\chapter{Lösungsansatz}
\label{sec:S4_Lösungsansatz}

\section{Mögliche Lösungen}

\subsection*{Gamification}

Im Hinblick auf die Problemstellung in Kapitel \ref{ch2:Problemstellunug}, sowie der in der Literatur beschriebenen Ansätze in Kapitel \ref{ch3:s:Gamification}, ergeben sich mehrere Möglichkeiten für eine Lösung. Ziel ist es die Kunden zum einen zu motivieren neue Geschäfte aufzusuchen und auf der anderen Seite eine Bindung zu bestehenden Geschäften zu erreichen. Da es sich hierbei um um einen Prozess handelt, der für den Kunden im ersten Schritt nur einen geringen Benefit darstellt, ist dieser ideal für Gamification geeignet \cite{Leigh.2012}. Ziel ist es somit die Besuche des Spielers durch eine Gamification interessant zu machen. Dies kann durch die verschiedenen Aspekte wie das Vergeben von Punkten für einen Einkauf (siehe Miles \& More \cite{Wagner.2005}, Payback \cite{Roesl.2005}) erreicht werden. Allerdings sind reine, einfache Punkteprogramme weit verbreitet und deren Nutzen umstritten \cite{Schmitt.2001}.
Daher ist es notwendig, das komplette Spektrum der Gamification zu betrachten.
Über das reine Punkte System hinaus, mandatetimeifestiert sich in der Literatur die Grundmenge "Points, Badges und Leaderboards".
Konkret bedeutet dies, neben der einfachen Ansammlung von Punkten, gibt es darüber hinaus Badges/Auszeichnungen welche das Engagement des einzelnen Nutzers widerspiegeln. Ein weiterer Aspekt sind die sogenannten Leaderboards bzw. Bestenlisten. Über diese wird wird ein Ansporn unter den jeweiligen Spielern erzeugt die anderen Spieler zu übertrumpfen. Vergleicht man diese Ansätze mit denen der typischen Spieler Profile nach \citep{Bartle.2004} so erkennt man, dass nicht alle dieser mit den Standard Elementen bedient werden.

\begin{figure}[H]
\begin{center}
\includegraphics[width=120mm]{images/ch4_img01_playertypes.png}
\caption{Spielertypen nach \textcite{Bartle.2004}}
\label{img:ch03_img02_flow}
\end{center}
\end{figure}

Die rudimentären Elemente der Gamification richten sich vor allem an Achievers und Killers.
Die Frage stellt sich hierbei, wie Socializers und Explorers erreicht werden können.
Hierzu kann das von \cite{Zichermann.2011} in Kapitel \ref{ch3:s:Gamification} erweiterte Modell SAPS genutzt werden, welches auch auf die Interaktion unter den Menschen eingeht. Im konkreten Beispiel schlägt Zichermann vor, den Spielern die Möglichkeit zu geben bestimmte Positionen innerhalb eines Spiels einzunehmen. Durch diese sollen einzelne Spieler eine höhere Position erhalten als andere und damit auch eine Interaktion untereinander fördern. Damit kann können zumindest zu einem gewissen Teil die Socializers abgedeckt werden.
Es ist jedoch zu beachten, dass für die vollständige Abdeckung aller Spielertypen konkret eine Spielmechanik erarbeitet werden muss, welche die übrigen zwei Spielertypen explizit bedient. Möglichkeiten sind hierbei eine Interaktion zwischen den Spielern zu schaffen, stellen z.B. das Erstellen von Teams/Gilden dar, welche gemeinsam aktiv handeln. Gemeinsame Ziele können so einfacher verfolgt werden aber auch speziell durch das Spiel ein Zusammenhalt explizit gefördert werden. Dies kann z.B. anhand von Bonuspunkten oder aber von Handlungen die explizit die Interaktion von mehreren Spielern erfordern. Ein Beispiel hierfür ist das Spiel Ingress von Google welches für die Erstellung von Portalen im späteren Spiel umso mehr Spieler benötigt, je höher dessen Level sein soll. \cite{Celino.2013}
Andere Ansätze sind kooperative Quests, Handel von Spielitems, sowie das Fördern von Kommunikation zwischen den einzelnen Spielern. \cite{Caon.2013} Ein weiterer Aspekt sind neue Bekanntschaften, die durch ortsbezogene Spiele gemacht werden. In einer Umfrage von 6400 Ingress Spielern gaben 84\% an neue Leute kennen gelernt zu haben und zählten dies zu einem wichtigen Faktor des Spiels \cite{Ingress.2014}.
Explorer untersuchen Ihre Umgebung und möchten neue Dinge entdecken. Sie fragen nach zusätzlichen Informationen. Für diese Spieler ist es z.B. interssant n neue unbekannte Orte innerhalb des Spiels zu gelangen.

Dass die einzelnen Spielertypen nicht zu vernachlässigen sind wird in \cite{Bromley.2013} deutlich. Hierbei stellt der Autor am Beispiel von GoPets fest, dass durch die Orientierung an konkrete Spielertypen ein beachtliche Umsatzsteigerung erzielt werden kann. \cite{Bethke.2007} Allerdings ist es wichtig für den Erfolg eines Spieles, dass alle Spielertypen adressiert werden.
\\\\
Nachdem die (Spiel-)Elemente der Gamification selbst identifiziert wurden, muss untersucht werden, wie die Interaktion der Spieler mit dem lokalen Einzelhandel aussehen kann.
Der Unterschied des Einzelhandels im Vergleich zum Onlinehandel ist der geografische bzw. physikalische Besuch. Der Kunde muss physikalisch vor Ort sein, um eine Dienstleistung wahrzunehmen oder aber ein Produkt zu kaufen. Daher ist es für den jeweiligen Händler von besonderem Interesse, dass ein potentieller Kunde sein Geschäft betritt.
Ziel ist es somit den Spieler durch die Spielmechanik in das jeweilige Geschäft zu bringen.
Hierbei gibt es zwei Mögliche Ansätze für ein ortsbezogenes Spiel. In der ersten und einfachsten Möglichkeit tritt das Geschäft selbst als festes Spielelement auf, dass direkt zum Spielerfolg beiträgt. In diesem Fall repräsentiert das Geschäft zum Beispiel ein TicTacToe Feld. Eine andere Möglichkeit ist es die Spieler indirekt zu unterstützten. In diesem Fall tritt das Geschäft z.B. als virtueller laden im Spiel auf, über den die Spieler entsprechende Items kaufen können.
Welche der beiden Möglichkeiten in Frage kommt, hängt von der Intention des einzelnen Händlers ab. Steht der Fokus des Spiels auf einem reinen Kontakt mit dem Spielelement, so 
muss sicher gestellt werden, dass die Spieler nicht einfach nur den jeweiligen Laden einfach nur betreten und dann wieder direkt verlassen. Im schlimmsten Fall kann dies sogar dazu führen, dass normale Kunden die nicht am Spiel selbst teilnehmen sich durch die Spieler gestört fühlen.
Da es hierfür keinerlei Erwartungswerte gibt, ist der Aspekt schwer abzuschätzen, allerdings wäre eine mögliche Gegenmaßenahme, die Verweildauer der Spieler zu erhöhen, in dem Sie z.B. ein kleines Rätsel lösen müssen oder für eine gewisse Zeit im entsprechenden Raum aufhalten müssen.
Des weiteren kann dem Händler auch die Möglichkeit gegeben werden, wiederkehrende Spieler bevorzugt zu behandeln. Ein Konzept welche z.B. bei Foursquare verwendet wird ist der Mayor einer Location. Der Mayor stellt dabei die Person,da welche die Lokalität innerhalb der letzten 60 Tage am häufigsten besucht hat. Diesem kann ein Händler über Foursquare besonderes Angebot machen, wie zum Beispiel ein Kaffee zum halben Preis.\cite{Lindqvist.2011} Ähnliche Konzepte können hier analog ebenfalls verwenden werden.
\\\\
Da das Hauptaugenmerk darauf liegt, dass der Spieler mit dem jeweiligen Geschäft interagiert, muss sichergestellt werden, dass dieser physikalisch vor Ort ist. 
Zunächst kann darauf

Das Betrügen in \cite{He.2011}


%
\subsection*{Anforderungen an ein Geogameframework}

\subsection*{Relokalisierbarkeit von ortsbezogenen Spielen}

\subsection*{Freie Geobasisdaten und Möglichkeiten der kommerziellen Nutzung}

\section{Gewählter Lösungsansatz}


